\chapter{Discussion and Interpretability}
\label{ch:discussion}

This chapter analyzes the interpretability of our classical guitar RubricNet model, evaluates the monotonicity of the extracted descriptors, and discusses the pedagogical implications of the feature design.

\section{Descriptor Influence Analysis}
To understand which technical features drive the difficulty predictions, we examine the descriptor score ranges. For a given trained model and test set, the influence of descriptor $i$ is calculated as the difference between its maximum and minimum activated subnetwork outputs:
\begin{equation}
    \text{Influence}(x_i) = \max_{x \in \mathcal{D}_{\text{test}}} g_i(x_i) - \min_{x \in \mathcal{D}_{\text{test}}} g_i(x_i)
\end{equation}
A larger range indicates that the subnetwork for descriptor $i$ has a greater impact on shifting the final aggregated score $S(x)$ across class boundaries.

\subsection{Top Descriptors in V2}
In RubricNet V2 (trained on fold 0 of seed 0), the five most influential descriptors were:
\begin{enumerate}
    \item \texttt{log\_total\_notes} (range = 1.989)
    \item \texttt{avg\_string\_jump} (range = 1.972)
    \item \texttt{total\_notes} (range = 1.942)
    \item \texttt{fret\_entropy} (range = 1.898)
    \item \texttt{chord\_ratio} (range = 1.742)
\end{enumerate}
These results indicate that global piece length (\texttt{total\_notes}) and left-hand fretboard coverage variety (\texttt{fret\_entropy}) act as primary indicators of difficulty, alongside right-hand plucking precision demands (\texttt{avg\_string\_jump}).

\subsection{Top Descriptors in V3}
In RubricNet V3, the five most influential descriptors shifted to:
\begin{enumerate}
    \item \texttt{log\_total\_notes} (range = 1.931)
    \item \texttt{total\_notes} (range = 1.930)
    \item \texttt{shift\_rate} (range = 1.911)
    \item \texttt{open\_string\_ratio} (range = 1.852)
    \item \texttt{string\_entropy} (range = 1.833)
\end{enumerate}
Here, the frequency of vertical position shifts (\texttt{shift\_rate}) rises to the top, while the fraction of open strings (\texttt{open\_string\_ratio}) acts as a strong negative weight (as open strings lower execution difficulty).

\section{Validation of Monotonicity}
For a difficulty estimation rubric to be pedagogically useful, it must exhibit positive monotonicity: increasing a descriptor's value should generally increase the predicted difficulty. In standard neural networks, this relationship is non-monotonic due to high-dimensional interactions. 

In RubricNet, the per-descriptor subnetwork outputs $g_i(x_i)$ plotted against the true difficulty classes exhibit clear, steady monotonic trends. For example, as the difficulty class increases from 0 (Easy) to 7 (Hard), the subnetwork outputs for \texttt{fret\_entropy}, \texttt{shift\_rate}, and \texttt{barre\_ratio} rise consistently. This mathematical constraint ensures that if a student asks *why* a piece is graded as a Class 5 instead of a Class 3, the model can provide a transparent answer (e.g., "The piece requires 15\% more position shifts and a higher barre ratio").

\section{Context-Free vs. Context-Aware Physical Proxies}
A key hypothesis in our V3 feature design was that physical execution demands are highly dependent on time and tempo. Table~\ref{tab:proxy_comparison} summarizes our findings when comparing context-free descriptors with their context-aware equivalents.

\begin{table}[h]
\centering
\caption{Comparison of Physical Proxies}
\label{tab:proxy_comparison}
\begin{tabular}{L{4.5cm}L{4.5cm}L{5cm}}
\toprule
\textbf{Context-Free Feature} & \textbf{Context-Aware Feature} & \textbf{Pedagogical Findings} \\
\midrule
\texttt{avg\_chord\_stretch} \newline (fret distance) & \texttt{avg\_stretch\_velocity\_beats} \newline (stretch change per beat) & Wide stretches are much harder to prepare when notes are played in rapid succession. \\
\midrule
\texttt{avg\_position\_shift} \newline (fret shift size) & \texttt{avg\_position\_shift\_speed\_beats} \newline (fret shift distance per beat) & Large jumps along the neck are trivial during rests, but highly difficult over sixteenth notes. \\
\midrule
\texttt{arpeggio\_density} \newline (string changes) & \texttt{polyphonic\_arpeggio\_intensity} \newline (density $\times$ active polyphony) & Plucking across strings is harder when maintaining multiple contrapuntal voices. \\
\bottomrule
\end{tabular}
\end{table}

Under Spearman correlation auditing, the context-aware velocity features exhibited stronger correlations with true difficulty than their context-free counterparts. For instance, combining position shift distance with beat duration (\texttt{avg\_position\_shift\_speed\_beats}) resolved false classifications where slow, simple pieces with a single large jump were over-graded.

\section{Limitations}
Despite the high performance of RubricNet V3, several limitations remain:
\begin{enumerate}
    \item \textbf{Single-Voice Alignment Constraint}: Our Needleman-Wunsch rhythm-alignment pipeline collapses multi-voice classical guitar scores into a single sequence of note attacks for alignment. While sufficient for extracting global and left-hand features, it misses fine-grained voice-leading structures.
    \item \textbf{MIDI Dependency}: The alignment algorithm requires a matched MIDI file for each vector PDF score. Without a MIDI performance file, the PDF's rhythm cannot be parsed, limiting the system's ability to grade arbitrary vector PDFs without human-in-the-loop matching.
\end{enumerate}
